% Chapitre 8 : Conclusion
\chapter{Conclusion et Perspectives}
\label{chap:conclusion}

\section{Synthèse du Projet}

Ce projet a permis de développer un système complet d'analyse forensique assistée par intelligence artificielle, combinant plusieurs technologies de pointe pour offrir une solution innovante et pratique.

\subsection{Objectifs Atteints}

Les objectifs initiaux du projet ont été pleinement réalisés :

\begin{itemize}
    \item \textbf{Détection d'objets} : Implémentation réussie de YOLOv8 pour l'identification automatique d'éléments sur les scènes analysées, avec des temps de réponse inférieurs à 200ms.
    
    \item \textbf{Extraction de texte} : Intégration de Tesseract OCR permettant de récupérer les informations textuelles présentes sur les images (plaques, panneaux, documents).
    
    \item \textbf{Génération d'hypothèses} : Utilisation de modèles de langage (Ollama/OpenAI) pour produire des analyses forensiques contextualisées et pertinentes.
    
    \item \textbf{Interface immersive} : Développement d'un avatar 3D interactif avec synchronisation labiale, rendant l'interaction plus naturelle et engageante.
    
    \item \textbf{Architecture moderne} : Mise en place d'une architecture microservices conteneurisée, facilitant le déploiement et la scalabilité.
\end{itemize}

\section{Difficultés Rencontrées}

\subsection{Défis Techniques}

Plusieurs défis ont été relevés durant le développement :

\begin{itemize}
    \item \textbf{Synchronisation Audio/Visème} : L'alignement précis entre l'audio généré et les mouvements labiaux de l'avatar a nécessité plusieurs itérations pour obtenir un résultat naturel.
    
    \item \textbf{Latence du Pipeline} : L'enchaînement Détection → OCR → LLM introduit une latence cumulative, partiellement compensée par le streaming.
    
    \item \textbf{Contexte de Conversation} : Maintenir un historique de conversation pertinent sans dépasser les limites de tokens des LLM a demandé une conception soignée.
\end{itemize}

\subsection{Solutions Apportées}

\begin{table}[H]
\centering

\label{tab:solutions}
\begin{tabular}{@{}p{5cm}p{9cm}@{}}
\toprule
\textbf{Problème} & \textbf{Solution} \\
\midrule
Latence élevée & Streaming SSE + Parallélisation des services \\
Mémoire GPU limitée & Chargement paresseux des modèles (lazy loading) \\
Lip-sync imprécis & Mapping Rhubarb → RPM visèmes + interpolation \\
Historique trop long & Résumé automatique + sliding window \\
\bottomrule
\end{tabular}
\caption{Problèmes et solutions}
\end{table}

\section{Compétences Acquises}

Ce projet a permis de développer et renforcer de nombreuses compétences :

\subsection{Compétences Techniques}

\begin{itemize}
    \item Maîtrise de FastAPI et des architectures REST/SSE
    \item Intégration de modèles de Deep Learning (YOLO, LLM)
    \item Développement 3D avec Three.js et React Three Fiber

\end{itemize}

\subsection{Compétences Transversales}

\begin{itemize}
    \item Architecture logicielle et design patterns
    \item Gestion de projet et méthodologie Agile
    \item Documentation technique et rédaction
    \item Tests et assurance qualité
\end{itemize}

\section{Perspectives d'Amélioration}
\begin{enumerate}
    \item \textbf{Modèle YOLO Spécialisé} : Entraîner un modèle YOLOv8 sur un dataset forensique pour améliorer la détection d'objets spécifiques (armes, preuves, véhicules).
    
    \item \textbf{Support Multi-langues} : Étendre l'OCR et le LLM pour supporter le français et d'autres langues.
    
    \item \textbf{Authentification} : Ajouter un système d'authentification JWT pour sécuriser l'accès aux API.
    
    \item \textbf{Annotations Visuelles} : Afficher les bounding boxes directement sur les images dans l'interface.
\end{enumerate}

\section{Conclusion Générale}

Le projet AI Forensic Avatar démontre le potentiel de l'intelligence artificielle pour assister les professionnels de l'investigation forensique. En combinant vision par ordinateur, traitement du langage naturel et interfaces immersives, nous avons créé un outil qui peut :

\begin{itemize}
    \item Accélérer l'analyse préliminaire des scènes
    \item Fournir des hypothèses objectives et documentées
    \item Rendre l'IA plus accessible grâce à une interface conversationnelle naturelle
\end{itemize}

Ce projet illustre également l'importance d'une architecture logicielle bien pensée, permettant l'évolution et l'extension du système. Les technologies choisies (FastAPI, React, Three.js, Docker) constituent une base solide pour les développements futurs.

Enfin, ce travail ouvre la voie à de nombreuses améliorations et applications connexes, confirmant que l'IA peut être un allié précieux dans le domaine forensique, tout en restant un outil au service de l'expertise humaine.

\vspace{1cm}

\begin{flushright}
\textit{« L'intelligence artificielle ne remplace pas l'enquêteur,\\
elle augmente ses capacités d'observation et d'analyse. »}
\end{flushright}
